
\section{Quantitative fit tables}
\label{app:quant_fits}

This appendix records compact quantitative tables supporting Section~\ref{sec:quantitative_closure}. All numbers are reproducible from the scripts in Appendix~\ref{app:repro_code}.

\subsection{Gauss digit-law fit for continued fractions}
\label{app:quant_fits_gauss_digits}

Section~\ref{sec:gauss_measure_digit_law} records the Gauss invariant measure and the closed-form digit probabilities~\eqref{eq:gauss_digit_law}. As a numerical audit, we iterate the Gauss map and compare empirical digit frequencies (digits $1$ through $K$ and a tail bin) to the theoretical probabilities.

\begin{table}[H]
\centering
\begin{tabular}{rrrr}
\toprule
$N$ & $K$ & $\ell^1$ error & $\ell^\infty$ error \\
\midrule
$10^4$   & 10 & $2.558203\times 10^{-2}$ & $6.309404\times 10^{-3}$ \\
$10^5$   & 10 & $6.297209\times 10^{-3}$ & $1.010596\times 10^{-3}$ \\
$10^6$   & 10 & $2.090263\times 10^{-3}$ & $4.741179\times 10^{-4}$ \\
\bottomrule
\end{tabular}
\caption{Gauss digit-law fit: empirical digit frequencies for the Gauss map compared to the theoretical digit probabilities~\eqref{eq:gauss_digit_law}. Reproducible via Experiment~IV in Appendix~\ref{app:repro_code}.}
\label{tab:gauss_digit_fit}
\end{table}

\subsection{Roof-function mean fit: the scale-time constant}
\label{app:quant_fits_roof_mean}

Section~\ref{sec:modular_flow_renormalization} identifies the modular geodesic flow as a suspension over the Gauss map with roof function $r(\xi)=-2\log\xi$. Proposition~\ref{prop:gauss_roof_mean} gives the closed-form mean $\mathbb{E}_\mu[r]=\pi^2/(6\log 2)$. The following table audits the ergodic average $\frac{1}{N}\sum_{j=1}^{N} r(G^j\xi)$ after burn-in, aggregated over multiple independent orbits (distinct initial points). We report the mean and standard deviation across orbits, together with the maximum absolute error across orbits.

\begin{table}[H]
\centering
\resizebox{\textwidth}{!}{%
\begin{tabular}{rrrrrrr}
\toprule
$N$ & mean over orbits & std over orbits & theoretical mean & absolute error & relative error & max absolute error \\
\midrule
$10^4$   & $2.369719307356\times 10^{0}$ & $1.897402\times 10^{-2}$ & $2.373138220831\times 10^{0}$ & $3.418913\times 10^{-3}$ & $1.440672\times 10^{-3}$ & $3.843985\times 10^{-2}$ \\
$5\times 10^4$ & $2.373896316300\times 10^{0}$ & $6.455681\times 10^{-3}$ & $2.373138220831\times 10^{0}$ & $7.580955\times 10^{-4}$ & $3.194485\times 10^{-4}$ & $1.336913\times 10^{-2}$ \\
$10^5$   & $2.373651962619\times 10^{0}$ & $4.731725\times 10^{-3}$ & $2.373138220831\times 10^{0}$ & $5.137418\times 10^{-4}$ & $2.164820\times 10^{-4}$ & $9.059725\times 10^{-3}$ \\
$3\times 10^5$ & $2.372742816984\times 10^{0}$ & $3.024095\times 10^{-3}$ & $2.373138220831\times 10^{0}$ & $3.954038\times 10^{-4}$ & $1.666164\times 10^{-4}$ & $8.243719\times 10^{-3}$ \\
$10^6$   & $2.372510303811\times 10^{0}$ & $1.264399\times 10^{-3}$ & $2.373138220831\times 10^{0}$ & $6.279170\times 10^{-4}$ & $2.645935\times 10^{-4}$ & $2.444371\times 10^{-3}$ \\
\bottomrule
\end{tabular}
}
\caption{Roof-function mean fit: empirical averages of $r(G^j\xi)$ compared to $\pi^2/(6\log 2)$, aggregated over multiple independent orbits (distinct initial points). Reproducible via Experiment~VII in Appendix~\ref{app:repro_code}.}
\label{tab:roof_mean_fit}
\end{table}

\subsection{Slice--sampling coefficient recovery: an explicit bound for \texorpdfstring{$E_4$}{E4}}
\label{app:quant_fits_slice_sampling_e4}

Section~\ref{sec:cusp_qexp_slice_sampling} gives a slice--sampling estimator~\eqref{eq:scan_estimator} and the discrepancy bound~\eqref{eq:finiteN_coeff_bound}. For the Eisenstein series
\begin{equation}
E_4(\tau)=1+240\sum_{n\ge 1}\sigma_3(n)\,q^n,
\end{equation}
we can bound $\Var(g_{n,y})$ explicitly (Appendix~\ref{app:variation_bound_e4_explicit}), producing a fully numerical, constant-explicit error upper bound. The following table reports the observed coefficient-recovery error for $a_1=240$ at $y=1$ (golden branch scan), together with the explicit bound and the ratio.

\begin{table}[H]
\centering
\resizebox{\textwidth}{!}{%
\begin{tabular}{rrrrrr}
\toprule
$N$ & $D_N^\ast$ & $|\widehat a_{1,N}(1)-240|$ & bound & $\log_{10}$ trunc bound & ratio \\
\midrule
100   & $1.731962\times 10^{-2}$ & $3.304239\times 10^{0}$ & $1.224815\times 10^{2}$ & $-2171.830$ & $2.697745\times 10^{-2}$ \\
200   & $8.945740\times 10^{-3}$ & $2.714377\times 10^{0}$ & $6.326283\times 10^{1}$ & $-2171.830$ & $4.290634\times 10^{-2}$ \\
500   & $3.598834\times 10^{-3}$ & $6.104705\times 10^{-2}$ & $2.545037\times 10^{1}$ & $-2171.830$ & $2.398670\times 10^{-3}$ \\
1000  & $1.759034\times 10^{-3}$ & $6.100095\times 10^{-2}$ & $1.243960\times 10^{1}$ & $-2171.830$ & $4.903769\times 10^{-3}$ \\
2000  & $9.488411\times 10^{-4}$ & $6.074182\times 10^{-2}$ & $6.710051\times 10^{0}$ & $-2171.830$ & $9.052363\times 10^{-3}$ \\
5000  & $4.415762\times 10^{-4}$ & $5.861210\times 10^{-2}$ & $3.122756\times 10^{0}$ & $-2171.830$ & $1.876935\times 10^{-2}$ \\
10000 & $3.010556\times 10^{-4}$ & $5.070685\times 10^{-2}$ & $2.129017\times 10^{0}$ & $-2171.830$ & $2.381702\times 10^{-2}$ \\
\bottomrule
\end{tabular}
}
\caption{Slice--sampling recovery of the $E_4$ coefficient $a_1=240$ at $y=1$ using a golden-branch scan orbit. The error bound is the explicit instance of Theorem~\ref{thm:finiteN_coeff_bound} with a closed-form $\Var(g_{1,1})$ bound (Appendix~\ref{app:variation_bound_e4_explicit}). The truncation-bound column reports $\log_{10}$ of the truncation contribution to the recovered coefficient bound (Appendix~\ref{app:eisenstein_tail_bounds}), showing it is negligible at the chosen truncation depth. Reproducible via Experiment~V in Appendix~\ref{app:repro_code}.}
\label{tab:e4_slice_sampling_bound}
\end{table}

\subsection{Scale trade-off in height \texorpdfstring{$y$}{y}: constant-explicit bound versus observed error}
\label{app:quant_fits_slice_sampling_y_tradeoff}

Section~\ref{sec:scale_tradeoffs} emphasizes that increasing height $y$ both smooths the slice function (suppressing high-$q$ modes) and amplifies coefficient recovery by the factor $\e^{2\pi n y}$. The net effect is not monotone in $y$. The following table audits this trade-off for recovering the $E_4$ coefficient $a_1=240$ at fixed sample size $N=5000$ (golden branch scan), using the explicit $\Var(g_{1,y})$ bound from Appendix~\ref{app:variation_bound_e4_explicit}. Truncation error is negligible at the chosen truncation depth and heights (see Appendix~\ref{app:eisenstein_tail_bounds}).

\begin{table}[H]
\centering
\resizebox{\textwidth}{!}{%
\begin{tabular}{rrrrrr}
\toprule
$y$ & $|q|=\e^{-2\pi y}$ & $|\widehat a_{1,N}(y)-240|$ & bound & $\log_{10}$ trunc bound & ratio \\
\midrule
0.4 & $8.100\times 10^{-2}$ & $1.490629\times 10^{-2}$ & $3.925313\times 10^{0}$ & $-861.993$ & $3.797479\times 10^{-3}$ \\
0.6 & $2.305\times 10^{-2}$ & $7.979010\times 10^{-3}$ & $2.213339\times 10^{0}$ & $-1298.620$ & $3.604966\times 10^{-3}$ \\
0.8 & $6.561\times 10^{-3}$ & $1.719322\times 10^{-2}$ & $2.153547\times 10^{0}$ & $-1735.228$ & $7.983676\times 10^{-3}$ \\
1.0 & $1.867\times 10^{-3}$ & $5.861210\times 10^{-2}$ & $3.122756\times 10^{0}$ & $-2171.830$ & $1.876935\times 10^{-2}$ \\
1.2 & $5.315\times 10^{-4}$ & $2.054972\times 10^{-1}$ & $6.831308\times 10^{0}$ & $-2608.431$ & $3.008168\times 10^{-2}$ \\
\bottomrule
\end{tabular}
}
\caption{Height trade-off for slice--sampling recovery of the $E_4$ coefficient $a_1=240$ at fixed $N=5000$ (golden branch scan). The bound is the explicit instance of Theorem~\ref{thm:finiteN_coeff_bound} using Appendix~\ref{app:variation_bound_e4_explicit}. The truncation-bound column reports $\log_{10}$ of the truncation contribution to the recovered coefficient bound, showing it is astronomically negligible at the chosen truncation depth. Reproducible via Experiment~VI in Appendix~\ref{app:repro_code}.}
\label{tab:e4_slice_sampling_y_tradeoff}
\end{table}

\subsection{Certified \texorpdfstring{$S$}{S}-invariance check for \texorpdfstring{$j$}{j}: truncation bound versus observed difference}
\label{app:quant_fits_j_certified}

While Experiment~III provides a high-precision numerical sanity check of $j(\tau)\approx j(-1/\tau)$, a reviewer may ask for a certified truncation-error budget. Experiment~VIII computes $j$ from truncated $q$-series for $E_4,E_6$ and propagates explicit tail bounds (Appendix~\ref{app:eisenstein_tail_bounds}) through the algebraic expression $j=1728 E_4^3/(E_4^3-E_6^2)$ to produce a certified bound $B(\tau)$ such that $|j(\tau)-j^{(N)}(\tau)|\le B(\tau)$. Since $j(\tau)=j(-1/\tau)$ exactly, we obtain the certified inequality
\begin{equation}
|j^{(N)}(\tau)-j^{(N)}(-1/\tau)|\le B(\tau)+B(-1/\tau).
\end{equation}

\begin{table}[H]
\centering
\begin{tabular}{rrrr}
\toprule
$N$ & $|j^{(N)}(\tau)-j^{(N)}(-1/\tau)|$ & $B(\tau)+B(-1/\tau)$ & ratio \\
\midrule
20 & $3.905295\times 10^{0}$ & $7.350896\times 10^{0}$ & $5.312680\times 10^{-1}$ \\
30 & $9.581200\times 10^{-5}$ & $1.716618\times 10^{-4}$ & $5.581441\times 10^{-1}$ \\
40 & $1.398789\times 10^{-9}$ & $2.372240\times 10^{-9}$ & $5.896490\times 10^{-1}$ \\
\bottomrule
\end{tabular}
\caption{Certified $S$-invariance check for the $j$-invariant at $\tau=0.3+0.2\iu$ using $q$-series truncation depth $N$. Reproducible via Experiment~VIII in Appendix~\ref{app:repro_code}.}
\label{tab:j_certified_S}
\end{table}

\subsection{Slice--sampling coefficient recovery: an explicit bound for \texorpdfstring{$E_6$}{E6}}
\label{app:quant_fits_slice_sampling_e6}

To show that the slice--sampling pipeline is not specific to $E_4$, we also audit the recovery of the first Fourier coefficient of the weight-$6$ Eisenstein series
\begin{equation}
E_6(\tau)=1-504\sum_{n\ge 1}\sigma_5(n)\,q^n,
\end{equation}
whose first coefficient is $a_1=-504$. Using $\sigma_5(n)\le \zeta(5)n^5$ and the same Fourier-coefficient variation bound strategy as Appendix~\ref{app:variation_bound_e4_explicit} (implemented explicitly in the reproduction script), we obtain a constant-explicit bound for $\Var(g_{1,1})$ and hence for the finite-$N$ recovery error. The following table is reproducible via Experiment~IX.

\begin{table}[H]
\centering
\resizebox{\textwidth}{!}{%
\begin{tabular}{rrrrrr}
\toprule
$N$ & $D_N^\ast$ & $|\widehat a_{1,N}(1)-(-504)|$ & bound & $\log_{10}$ trunc bound & ratio \\
\midrule
100   & $1.731962\times 10^{-2}$ & $3.507447\times 10^{0}$ & $1.824075\times 10^{2}$ & $-1620.638$ & $1.922863\times 10^{-2}$ \\
200   & $8.945740\times 10^{-3}$ & $2.704241\times 10^{0}$ & $9.421512\times 10^{1}$ & $-1620.638$ & $2.870284\times 10^{-2}$ \\
500   & $3.598834\times 10^{-3}$ & $6.341491\times 10^{-2}$ & $3.790235\times 10^{1}$ & $-1620.638$ & $1.673113\times 10^{-3}$ \\
1000  & $1.759034\times 10^{-3}$ & $6.301880\times 10^{-2}$ & $1.852587\times 10^{1}$ & $-1620.638$ & $3.401665\times 10^{-3}$ \\
2000  & $9.488411\times 10^{-4}$ & $6.200086\times 10^{-2}$ & $9.993045\times 10^{0}$ & $-1620.638$ & $6.204401\times 10^{-3}$ \\
5000  & $4.415762\times 10^{-4}$ & $5.749970\times 10^{-2}$ & $4.650611\times 10^{0}$ & $-1620.638$ & $1.236390\times 10^{-2}$ \\
10000 & $3.010556\times 10^{-4}$ & $4.752859\times 10^{-2}$ & $3.170671\times 10^{0}$ & $-1620.638$ & $1.499007\times 10^{-2}$ \\
\bottomrule
\end{tabular}
}
\caption{Slice--sampling recovery of the $E_6$ coefficient $a_1=-504$ at $y=1$ using a golden-branch scan orbit. The bound is the explicit instance of Theorem~\ref{thm:finiteN_coeff_bound} with the closed-form $\Var(g_{1,1})$ bound from Appendix~\ref{app:variation_bound_e6_explicit}. The truncation-bound column reports $\log_{10}$ of the truncation contribution to the recovered coefficient bound (Appendix~\ref{app:eisenstein_tail_bounds}). Reproducible via Experiment~IX in Appendix~\ref{app:repro_code}.}
\label{tab:e6_slice_sampling_bound}
\end{table}

\subsection{Slice--sampling closure on a cusp form: \texorpdfstring{$\Delta$}{Delta} and Ramanujan \texorpdfstring{$\tau(n)$}{tau(n)}}
\label{app:quant_fits_slice_sampling_delta}

To fully close the stairway chain on a single arithmetic object, we apply the slice--sampling estimator~\eqref{eq:scan_estimator} to the weight-$12$ cusp form
\begin{equation}
\Delta(\tau)=\sum_{n\ge 1}\tau(n)\,q^n,
\end{equation}
recovering $\tau(n)$ from scan-orbit samples of the slice function $x\mapsto \Delta(x+\iu y)$.
We use the same discrepancy-controlled bound (Theorem~\ref{thm:finiteN_coeff_bound}) and evaluate $\Delta$ via its truncated $q$-series with a certified tail majorant based on Proposition~\ref{prop:tau_global_bound}.

\begin{table}[H]
\centering
\begin{tabular}{rrrrrr}
\toprule
$n$ & $\tau(n)$ & $|\widehat\tau_{n,N}(y)-\tau(n)|$ & bound & $\log_{10}$ trunc bound & ratio \\
\midrule
1 & 1 & $4.836469\times 10^{-4}$ & $2.380285\times 10^{-2}$ & $-204.142$ & $2.031886\times 10^{-2}$ \\
2 & $-24$ & $4.407156\times 10^{-3}$ & $1.700035\times 10^{-1}$ & $-203.051$ & $2.592392\times 10^{-2}$ \\
4 & $-1472$ & $9.061871\times 10^{-1}$ & $4.601548\times 10^{1}$ & $-200.868$ & $1.969309\times 10^{-2}$ \\
\bottomrule
\end{tabular}
\caption{Slice--sampling recovery of $\tau(n)$ from scan-orbit samples of $\Delta(x+\iu y)$ on the golden branch ($y=0.4$, $N=5000$). The bound is the explicit instance of Theorem~\ref{thm:finiteN_coeff_bound} with a certified $\Var(g_{n,y})$ upper bound and a certified truncation-tail budget based on Proposition~\ref{prop:tau_global_bound}. Reproducible via Experiment~X in Appendix~\ref{app:repro_code}.}
\label{tab:delta_tau_slice_sampling}
\end{table}

\subsection{Certified Hecke closure from recovered coefficients}
\label{app:quant_fits_tau_hecke_closure}

For a normalized weight-$12$ eigenform one has the prime-square recursion
\begin{equation}
\tau(p^2)=\tau(p)^2-p^{11}\tau(1).
\end{equation}
Using the recovered values $\widehat\tau(1),\widehat\tau(2),\widehat\tau(4)$ and their certified bounds from Table~\ref{tab:delta_tau_slice_sampling}, we obtain a certified inequality of the form
\begin{equation}
\left|\widehat\tau(4)-\big(\widehat\tau(2)^2-2^{11}\widehat\tau(1)\big)\right|\le \text{budget},
\end{equation}
where the right-hand side is obtained by explicit error propagation.

\begin{table}[H]
\centering
\begin{tabular}{rrr}
\toprule
lhs & budget & ratio \\
\midrule
$5.599303\times 10^{-1}$ & $1.029232\times 10^{2}$ & $5.440272\times 10^{-3}$ \\
\bottomrule
\end{tabular}
\caption{Certified Hecke closure for $p=2$ using scan-recovered coefficients on the golden branch ($y=0.4$, $N=5000$). Reproducible via Experiment~X.}
\label{tab:tau_hecke_certified_closure}
\end{table}


